\chapter{Background and Related Work}
\label{chap:background}

This chapter develops the concepts needed to compare MI between two
categorical populations. It begins with entropy, pointwise MI, and population
MI\@. It then reviews plug-in estimation, first-order inference, the
Welch--Satterthwaite approximation, and prior comparisons of information
quantities. The final section identifies the gap addressed by the expanded
method.

\section{Information-theoretic foundations}

\subsection{Entropy and information content}

For a discrete random variable \(X\) with probability mass function
\(p_X(x)\), the information content of outcome \(x\) is
\begin{equation}
  h_X(x)=-\log p_X(x).
\end{equation}
Rare outcomes have larger information content because they are less expected.
Using natural logarithms gives information in nats. The Shannon entropy is the
expected information content,
\begin{equation}
  H(X)=-\sum_x p_X(x)\log p_X(x).
  \label{eq:entropy}
\end{equation}
Entropy summarises the uncertainty of one random variable before its value is
observed \citep{shannon1948}.

For a pair \((X,Y)\), the joint entropy \(H(X,Y)\) is defined from the joint
probabilities in the same way. The conditional entropy
\(
H(X\mid Y)=H(X,Y)-H(Y)
\)
measures the uncertainty remaining in \(X\) after observing \(Y\).

\subsection{Pointwise and average mutual information}

Let \(P=(p_{ij})\) be the joint distribution of categorical variables
\(X\in\{1,\ldots,r\}\) and \(Y\in\{1,\ldots,c\}\). Write
\begin{equation}
  p_{i+}=\sum_{j=1}^{c}p_{ij},
  \qquad
  p_{+j}=\sum_{i=1}^{r}p_{ij}
\end{equation}
for the row and column margins. The pointwise mutual information associated
with cell \((i,j)\) is
\begin{equation}
  \pmi_P(i,j)
  =\log\left(\frac{p_{ij}}{p_{i+}p_{+j}}\right).
  \label{eq:background-pmi}
\end{equation}
The denominator is the cell probability under independence with the same
margins. A positive value means that the cell occurs more often than
independence predicts; a negative value means that it occurs less often.

Mutual information is the probability-weighted average of these pointwise
values:
\begin{equation}
  I(P)
  =\sum_{i=1}^{r}\sum_{j=1}^{c}p_{ij}\pmi_P(i,j).
  \label{eq:background-mi}
\end{equation}
Equivalently,
\begin{equation}
  I(X;Y)=H(X)+H(Y)-H(X,Y).
  \label{eq:mi-entropy-form}
\end{equation}
The quantity is nonnegative and equals zero exactly when
\(p_{ij}=p_{i+}p_{+j}\) for every cell, including zero-probability cells.

Equation~\eqref{eq:background-mi} is also the Kullback--Leibler divergence
between the joint distribution and the product of its margins. MI therefore
measures the overall departure from independence, not one particular
shape of association.

This scalar summary deliberately discards much of the information in the
joint distribution. Two populations can have the same MI while assigning
probability to different cells, having different margins, or exhibiting
different association patterns. Equality of MI is therefore weaker than
equality of distributions. This distinction is central to the present test:
the null fixes one scalar functional while leaving the remaining features of
each population unrestricted.

The logarithm base determines the unit but not the inferential problem. If
base two is used, every pointwise value and MI value is divided by
\(\log 2\). The standard error is divided by the same constant, so a
standardised statistic is unchanged. This thesis uses natural logarithms and
reports MI in nats.

\section{Estimating discrete mutual information}
\label{sec:background-mi-bias}

Suppose a sample of size \(n\) produces counts \(N_{ij}\). The empirical cell
probabilities are
\begin{equation}
  \widehat p_{ij}=\frac{N_{ij}}{n}.
\end{equation}
Substitution into Equation~\eqref{eq:background-mi} gives the plug-in estimator
\begin{equation}
  \Ihat(P)
  =\sum_{i,j:\widehat p_{ij}>0}
  \widehat p_{ij}
  \log\left(
  \frac{\widehat p_{ij}}
       {\widehat p_{i+}\widehat p_{+j}}
  \right).
  \label{eq:plugin-mi-background}
\end{equation}
The zero-cell convention follows from
\(\lim_{x\downarrow0}x\log x=0\). Empty cells contribute no direct term, but
they still affect the empirical margins and the finite-sample behaviour of the
estimator.

Under multinomial sampling, the entire table is random subject to its fixed
total. Cell counts are dependent because increasing one count leaves fewer
observations for the others. The plug-in estimator remains consistent for a
fixed alphabet, but its finite-sample distribution depends on the complete
probability table. Equal alphabet sizes and sample sizes do not imply equal
sampling behaviour when one population has much rarer cells.

The plug-in estimator is upward biased. Finite sampling makes an empirical
joint table appear more structured than its population distribution, which
increases the estimated departure from independence. For a fixed positive
\(r\times c\) support, the leading bias is
\begin{equation}
  \Bias\{\Ihat(P)\}
  =\frac{(r-1)(c-1)}{2n}+O(n^{-2}).
  \label{eq:mi-leading-bias}
\end{equation}
This result follows from the corresponding leading entropy biases and is
closely related to Miller's correction \citep{miller1955}; the same
histogram-MI term appears in \citet{moddemeijer1989}.
More complete analyses show that sparse sampling and unknown support can
produce substantial higher-order error beyond Equation
\eqref{eq:mi-leading-bias} \citep{paninski2003}.

The leading term can be seen from the entropy identity in Equation
\eqref{eq:mi-entropy-form}. For a discrete variable with \(m\) positive
categories, the plug-in entropy has leading bias \(-(m-1)/(2n)\). Applying
this result to \(H(X)+H(Y)-H(X,Y)\) gives
\begin{align}
  \Bias\{\Ihat(P)\}
  &\approx
  -\frac{r-1}{2n}
  -\frac{c-1}{2n}
  +\frac{rc-1}{2n}\\
  &=\frac{(r-1)(c-1)}{2n}.
  \label{eq:mi-bias-from-entropies}
\end{align}
The correction is positive because estimating the larger joint distribution
introduces more entropy bias than estimating the two margins.

Bias matters directly when comparing two populations. If
\(I(P)=I(Q)\) but \(n_P\ne n_Q\), the uncorrected difference has leading bias
\begin{equation}
  \frac{(r-1)(c-1)}{2}
  \left(\frac{1}{n_P}-\frac{1}{n_Q}\right).
\end{equation}
Removing the leading term is therefore part of the effect estimate, not
an optional refinement of its reference distribution.

\section{First-order inference for nonlinear functionals}

MI is a nonlinear function of the cell probabilities. When all population
cells are positive and the empirical table is close to its population value,
a first-order Taylor expansion replaces this nonlinear function locally by a
linear combination of the cell-probability errors. The resulting expansion is
often expressed through an influence function. The classical influence-curve
and functional delta-method treatments are given by \citet{hampel1974} and
\citet[Chapter~20]{vanderVaart1998}; \citet{kandasamy2015} apply related von
Mises methods to information functionals.

For MI, the observation-level influence function is centred pointwise MI:
\begin{equation}
  \psi_P(i,j)=\pmi_P(i,j)-I(P).
  \label{eq:mi-if-background}
\end{equation}
It has population mean zero. Its variance is
\begin{equation}
  V(P)
  =\Var_P\{\psi_P(X,Y)\}
  =\sum_{i,j}p_{ij}\{\pmi_P(i,j)-I(P)\}^2.
  \label{eq:mi-variance-background}
\end{equation}
The plug-in MI error is, to first order, the sample mean of these influence
values. Consequently,
\begin{equation}
  \Var\{\Ihat(P)\}\approx\frac{V(P)}{n}.
  \label{eq:first-order-mi-var-background}
\end{equation}
This leading variance for histogram MI was calculated by
\citet{moddemeijer1989}, who noted that the calculation also applies to
discrete systems. \citet{brillinger2004} gives the discrete non-null formula
explicitly, attributes that result to Moddemeijer, and states alongside it the
\(\chi^2/(2n)\) limit at independence. For temporally dependent observation
pairs, \citet{moddemeijer1999}
develops an estimator of the same sampling variance using methods that account
for dependence between observations. The present work assumes independent
observations within each sample and does not claim the first-order MI variance
formula as new.

This representation has two roles. First, it makes the central limit theorem
available because the leading estimation error is an average of independent,
mean-zero contributions. Second, it identifies the observation-level
quantity whose variance determines the standard error. The raw categories
themselves do not have a numerical variance relevant to MI; their effects on
the MI functional are represented by centred pointwise MI\@.

The expansion is local. It describes sampling fluctuations near a fixed
positive-support population and does not assert that the plug-in estimator is
normally distributed for every finite table. Its accuracy depends on the
sample size relative to the probability structure, not only on the total
number of observations.

The central limit theorem then gives a normal approximation when
\(V(P)>0\). At exact independence, every pointwise-MI value is zero and
\(V(P)=0\). The first-order term disappears, so a second-order Taylor expansion is
required. This is why the chi-squared limit for testing independence and the
normal limit for comparing positive MI values are not competing descriptions
of the same regular case.

\section{Wald and Welch--Satterthwaite inference}

\subsection{Normal Wald inference}

A Wald test estimates a parameter difference, divides it by an estimated
standard error, and compares the result with its asymptotic normal reference.
For two independent MI estimates, the first-order variance contributions add:
\begin{equation}
  \Var\{\Ihat(P)-\Ihat(Q)\}
  \approx
  \frac{V(P)}{n_P}+
  \frac{V(Q)}{n_Q}.
  \label{eq:wald-variance-background}
\end{equation}
Replacing the population quantities by estimates produces the Normal Wald
test developed formally in Chapter~\ref{chap:baselines}.

Wald inference is asymptotically valid under regularity conditions, but its
finite-sample calibration depends on both the quality of the linear expansion
and the stability of the estimated standard error. A normal reference does not
separately represent the uncertainty of the two estimated variance
contributions.

\subsection{The Welch--Satterthwaite principle}

Welch's test compares two means when their population variances may differ
\citep{welch1947}. The denominator contains two independently estimated
variance contributions. Satterthwaite's approximation represents their sum by
a scaled chi-squared variable with matching first two moments
\citep{satterthwaite1946}. The resulting effective degrees of freedom are
\begin{equation}
  \nu
  =
  \frac{\left(\sum_i k_i s_i^2\right)^2}
       {\sum_i (k_i s_i^2)^2/\nu_i},
  \label{eq:general-ws-background}
\end{equation}
where \(s_i^2\) are variance estimates, \(k_i\) are positive weights, and
\(\nu_i\) are their component degrees of freedom.

The final statistic uses a Student distribution rather than a normal
distribution. Small effective degrees of freedom produce heavier tails,
reflecting greater uncertainty in the estimated denominator. As the component
variance estimates become stable, the effective degrees of freedom increase
and the Student reference approaches the normal reference.

For a single scalar restriction, normal and chi-squared Wald notation are
equivalent. If \(T\) is asymptotically standard normal, then \(T^2\) is
asymptotically \(\chi^2_1\). The signed form is retained here because the
parameter is the scalar difference \(I(P)-I(Q)\), its sign is meaningful, and
the Welch refinement replaces the normal reference for \(T\) by a Student
reference. The \((r-1)(c-1)\)-degree chi-squared law used for one-table
independence is a different second-order result.

The same principle extends beyond ordinary sample variances.
\citet{panwall2002} account for variability in sandwich variance estimates
using approximate Student and $F$ references. \citet{bell2002} use
Satterthwaite moment matching for regression variance estimates under
multi-stage sampling after showing that traditional linearisation estimators
can have substantial negative bias. \citet{kauermann2001} examine the
variability and efficiency of sandwich covariance estimation and propose an
adjustment for the resulting confidence-interval undercoverage. Related
adjustments and practical guidance are given by \citet{fay2001} and
\citet{imbens2016}. These are close
methodological precedents: estimating the variability of a variance estimator
and using it in a finite-sample reference are established strategies. The
present task is to derive the required inputs for the complete discrete MI
functional and assess whether the resulting reference is useful.

The direction of the finite-sample problem is nevertheless different here.
The cited regression corrections respond to underestimated uncertainty or
anti-conservative coverage, so a larger variance estimate or heavier-tailed
reference can move inference in the required direction. Near independence,
the MI calculation instead gives
\(\E\widehat V\approx V+d/n\): the plug-in variance is shifted upward. A
heavier-tailed Student reference then reinforces, rather than offsets, the
resulting conservatism.

\section{Prior MI inference and related comparisons}

MI uncertainty also has a broader literature than its leading normal
approximation. \citet{roulston1999} derives entropy and MI error estimates
using error propagation and examines them numerically. \citet{hutter2002}
and \citet{hutterzaffalon2005} derive moments of the posterior distribution
of MI under a Dirichlet model, including higher-order variance and shape
information. Their posterior uncertainty has a different conditioning from
the repeated-sampling distribution studied here. Their results establish
that higher-order MI uncertainty is not a new topic, but do not directly
validate the two-sample Student reference in this thesis.

Hutcheson's test is a direct information-theoretic Welch predecessor. It
compares two Shannon diversity estimates with an estimated variance and a
Student reference, using
\begin{equation}
  \nu_H=
  \frac{(s_{H,P}^2+s_{H,Q}^2)^2}
       {(s_{H,P}^2)^2/n_P+(s_{H,Q}^2)^2/n_Q}
  \label{eq:hutcheson-df}
\end{equation}
as its effective degrees of freedom \citep{hutcheson1970}. Here each of
\(s_{H,P}^2\) and \(s_{H,Q}^2\) is the estimated sampling variance of the
corresponding entropy estimator \(\widehat H\), and each variance component
is assigned its sample size, not \(n_i-1\). Applying the same assignment to the
two MI variance contributions gives a clear historical baseline. It does not,
however, derive how the plug-in MI variance itself fluctuates: MI depends on a
joint distribution and both margins, so moving mass into one cell changes
pointwise values elsewhere in its row and column. The expanded calculation
retains these changes.

The question of comparing MI values also has predecessors. In a segregation
application, \citet{mora2009} develop inference for an MI index and discuss
pairwise comparisons across groups and periods; the later peer-reviewed account
places that MI index within the wider segregation literature \citep{mora2011}.
They do not derive the
variance-sensitivity degrees of freedom used here, but the general aim of
comparing MI values is not new. \citet{guha2014} use MI for a different two-sample
question: whether two continuous distributions are identical. Their MI is
between a pooled observation and its sample label; their null is equality of
distributions, not equality of two within-population MI values. For discrete
data, \citet{drake2014} develop an MI-based \(k\)-sample test for the same
distributional-equality question.

The exact-independence boundary has likewise been studied separately.
\citet{marinescu2025} use a second-order delta-method argument for MI under
independence. More generally, their expansion combines a normal term with a
linear combination of chi-squared terms; the normal term vanishes at exact
independence. That setting explains why a first-order positive-MI comparison
must not be presented as a replacement for a one-table independence test and
suggests why second-order effects can remain relevant near that boundary.

Table~\ref{tab:prior-work-scope} distinguishes the relevant established
ideas by the problem each addresses. This is not a claim that no other MI
comparison has been proposed. It states which pieces are used here and why
the final comparison still needs empirical validation.

\begin{table}[htbp]
\centering
\caption{Selected prior work, grouped by the question it addresses.}
\label{tab:prior-work-scope}
\begin{tabularx}{\textwidth}{@{}>{\raggedright\arraybackslash}p{0.24\textwidth}X@{}}
\toprule
Topic & Established result and relation to this thesis\\
\midrule
MI estimation & \citet{miller1955} and \citet{paninski2003}
establish bias and sparse-sample limitations of plug-in information estimates.\\
MI variance & \citet{moddemeijer1989} and \citet{brillinger2004}
give predecessors for the first-order MI variance, while
\citet{moddemeijer1999} treats dependent observation pairs. This thesis instead
studies the sampling uncertainty of the plug-in variance estimate itself in
the independent two-sample problem.\\
Welch principle & \citet{satterthwaite1946} and \citet{welch1947}
provide moment-matching ideas; \citet{hutcheson1970} applies the same structure
to Shannon diversity using component degrees of freedom \(n_i\).\\
Other comparisons & \citet{mora2009,mora2011} study an MI-based segregation
index. They do not derive the variance-sensitivity degrees of freedom used here.\\
Different MI nulls & \citet{guha2014} and \citet{drake2014} test
distributional equality;
\citet{marinescu2025} study one-table independence. These are not the
two-population equal-MI null.\\
Weak-null tests & \citet{chung2013} explain why a raw pooled-label
permutation need not be valid when the compared parameters agree but
\(P\ne Q\); \citet{janssen1997} gives an earlier studentized permutation
result for heterogeneous two-sample nulls.\\
\bottomrule
\end{tabularx}
\end{table}

\section{Resampling and the weak null}

Bootstrap methods can approximate a statistic's repeated-sampling behaviour
by drawing from fitted empirical populations, at greater computational cost
than a direct analytic table scan \citep{efron1979}. Sparse empirical support
can still be a problem. A label permutation instead pools the observations and
reallocates their group labels; it targets a different symmetry assumption.

Permutation requires a separate distinction. Pooling two samples and
permuting their group labels is exact under the strong null \(P=Q\), because
the observations are exchangeable. Under the weak null \(I(P)=I(Q)\), the two
complete distributions may differ, so raw label permutation generally targets
the wrong null. Studentized permutation can recover asymptotic validity for
some weak-null parameter comparisons \citep{janssen1997,chung2013}, but it is not the
standard one-table MI shuffle used to test independence.

The final comparison uses Normal Wald and the proposed Expanded Welch method,
which share the same effect estimate and standard error. Resampling is relevant
future work, but is not an exact finite-sample solution for the unrestricted
weak null without additional assumptions.

\section{Literature synthesis and methodological gap}

The literature establishes all of the broad ingredients used in this thesis:
plug-in MI estimation and bias correction, first-order MI variance, Wald
comparison of smooth functionals, and Welch--Satterthwaite moment matching.
It also contains prior comparisons of entropy and MI-related indices.

No single ingredient is presented as new. Shannon established MI; Miller and
later authors developed finite-sample bias analysis; Moddemeijer and
Brillinger give the first-order MI sampling variance, including a
dependent-observation extension; and Welch and Satterthwaite provide the
general moment-matching structure. Hutcheson already applies that structure to
an information quantity using component degrees of freedom \(n_i\).
Mora and Ruiz-Castillo also show that inference about differences in an MI
index predates this study. The contribution claimed here is narrower: an
explicit observation-level sensitivity for the \emph{estimated MI variance},
its use in a plug-in Welch--Satterthwaite calculation, and a controlled
evaluation under distinct equal-MI multinomial populations. The literature
review does not establish that no earlier derivation of this exact
sensitivity exists, so first priority is not claimed.

The remaining task is specific to MI. Hutcheson's direct information-theoretic
predecessor assigns each component its sample size; an ordinary-Welch analogue
uses \(n-1\). Both are fixed assignments. By contrast, established sandwich
adjustments derive effective degrees of freedom from the estimated variability
of the variance estimator itself. The plug-in MI variance requires its own
version of that calculation because it is a nonlinear function of the same
empirical joint table used to calculate MI. A change in one observation
modifies the cell probabilities, margins, pointwise-MI values, and their
weighted variance.

The expanded method addresses this gap by differentiating the complete MI
variance functional. The resulting observation-level sensitivity determines
the sampling variance of the variance estimate and, through Satterthwaite
moment matching, an MI-specific component degree of freedom. The contribution
is this explicit functional calculation and its evaluation, within the
established approach to variance uncertainty described above.
Chapters~\ref{chap:baselines} and~\ref{chap:expanded} develop this construction.
